\chapter{The coefficient bound and its transfer}\label{ch-bound}

The exact identity~\eqref{budget7} is now converted into a lower bound for
the defect $D=(3+\sqrt5)/2-A-B$, which by Section~\ref{sec-conclusion} equals
$(3+\sqrt5)/2-2\mathcal M$ for the canonical law of a maximiser. Three
ingredients are combined: directional Hermite inequalities that pay for the
harmonic moments $H^2$ and $K^2$, a radial comparison that turns the slack
moment $P$ into a linear payment in $V$ and $J$, and an atomic term supplied
by the finite support theorem of Chapter~\ref{ch-support}. Each ingredient
leaves a nonnegative remainder, and the point of the construction is that all
of them are kept in one identity rather than discarded.

Section~\ref{sec-directional} carries out the directional step and produces
the intermediate certificate~\eqref{certificate-start} with coefficient $d_0$.
Section~\ref{sec-reserve} finds a reserve inside the quartic slack by means of
the cubic-free contact polynomial $z_*$, improving $d_0$ to $d_1$.
Section~\ref{sec-conclusion} verifies the radial hypothesis for an actual
maximiser, adds the positive-mass self-pair term, and reaches the final
constant $d$, equivalently $\mathcal M\le c$; it then transfers the bound to
all isotropic laws and to finite graphs, proving Theorem~\ref{main}, and
records the exact scope of the result.

\section{Directional and radial control}\label{sec-directional}
% Ported verbatim from paper/top3_joint_moment_bound.tex lines 957-1084.
% Only cross-reference wording was adapted; no mathematics was altered.
Throughout this section assume $0<r_-\le R\le r_+<\infty$ almost surely,
with $r_+/r_-\le3$, and work in the branch $D<1/20$.
Equation~\eqref{budget4} implies
\begin{equation}\label{smallbranch}
 e<1/20,\quad x>9/10,\quad V<21/104<2/9,\quad C>5/3,
 \qquad q^2=\frac{3-V}{9(1-e)}\le\frac{20}{57}.
\end{equation}
In particular $\mathbf m\ne0$. Put $p_i=\langle T_i,\Psi_i(v)\rangle
=\E[Rg_i(t)]$. The elementary Hermite identities
\begin{align}\label{Hermite}
 \left(\frac{53}{28}t^2-\frac94t+\frac{41}{70}\right)-g_4(t)
       &=(1-t)(t-1/2)^2(t+2)\ge0,\notag\\
 (2t^2-37t/20+1/4)-g_3(t)
       &=(1-t)(t-1/2)^2\ge0
\end{align}
hold on $[-1,1]$. Define
\[
 \gamma_4=9/4-53\sqrt3/42,\quad
 \ell_4=9/4-511\sqrt3/420,\quad
 \ell_3=37/20-11\sqrt3/12.
\]
Here $0<\gamma_4<1/15$ (square the positive rational comparisons).
Since $\E Rt=x$, $\E Rt^2=L+2x/\sqrt3-C/3$, the first identity gives
\[
 -p_4\ge\gamma_4x+19C/420-53L/28
       \ge\ell_4-e/15-19V/1400-53L/28.
\]
We used $x\ge1-e$ and $C\ge\sqrt3-3V/10$; the latter follows from
$V=(\sqrt3-C)(\sqrt3+C)$ and $C>5/3$.
The second identity gives
\[
 -p_3\ge(37/20-4/\sqrt3)x+5C/12-2L
       \ge\ell_3+9e/40-V/8-2L.
\]
Here $4/\sqrt3-37/20>9/20$ and $x\le1-e/2$.
Define the nonnegative differences in these last two inequalities by
\begin{align*}
 r_4&=-p_4-\ell_4+e/15+19V/1400+53L/28,\\
 r_3&=-p_3-\ell_3-9e/40+V/8+2L.
\end{align*}
The polynomial identities, rather than a fixed sphere assumption,
justify these directional estimates for the actual weighted moments.

Set $\eta=1/10$, $\theta=6/25$, and
$\ell=53\eta/28+2\theta=937/1400$. Define the tensor squares
\[
 {\cal R}_4=a_4\left\|T_4+\frac{\eta}{2a_4}\Psi_4(v)\right\|^2,\qquad
 {\cal R}_3=a_3\left\|T_3+\frac{\theta}{2a_3}\Psi_3(v)\right\|^2.
\]
Their exact expansions give
\[
 a_4K^2=-\eta p_4-\frac{2\eta^2}{35a_4}+{\cal R}_4,\qquad
 a_3H^2=-\theta p_3-\frac{\theta^2}{10a_3}+{\cal R}_3.
\]
Define
\[
 E=\frac{5183683}{5659500},\quad
 V_c=\frac{145858703}{581042000},\quad
 J_c=\frac{7989533}{14526050},\quad
 {\cal R}_0={\cal R}_3+{\cal R}_4+\eta r_4+\theta r_3\ge0.
\]
Substituting the directional identities and \eqref{payment} into
\eqref{budget7}, then writing $L=\overline L-r_{\rm reg}$, gives the exact identity
\begin{equation}\label{certificate-start}
 D=d_0+E e+V_c V+J_c J-\ell\overline L
             +{\cal R}_0+r_{\rm rad}+\ell r_{\rm reg}.
\end{equation}
Here
\[
 d_0=\delta_{\mathrm q}+\eta\ell_4+\theta\ell_3
                 -\frac{2\eta^2}{35a_4}-\frac{\theta^2}{10a_3},
\]
which simplifies to \eqref{constants}. The coefficient checks are
\[
 E=b_e-\eta/15+9\theta/40,\quad
 V_c=b_V-19\eta/1400-\theta/8+3/50,\quad J_c=b_J-27/50.
\]
All the remainders in this identity refer to the same probability law.

Put $B_c=V_c-5\ell/18$ and $C_c=J_c-\ell/4$. The exact rational values are
\[
 B_c=\frac{42565069}{653672250},\qquad
 C_c=\frac{88944217}{232416800}.
\]
For $u=\sqrt e$, $z=\sqrt V$, and $j=\sqrt J$, the remaining quadratic has
the explicit completion
\begin{align}\label{squares}
 E e+V_c V+J_c J-\ell\overline L
 &=\frac{13}{200}\left(z-\frac{937}{273}qu\right)^2
   +\frac{19}{50}\left(j-\frac{2811}{2128}qu\right)^2\notag\\
 &\quad+\left(\frac{183}{200}-\frac{725029049}{278853120}q^2\right)e\notag\\
 &\quad+\frac{10481}{11319000}e
        +\frac{305491}{2614689000}V+\frac{625833}{232416800}J.
\end{align}
Indeed the last three fractions are $E-183/200$, $B_c-13/200$,
and $C_c-19/50$. Also
\[
 \frac{725029049}{278853120}
  =\frac{7\ell}{4}+\frac{\ell^2}{9(13/200)}
                          +\frac{9\ell^2}{16(19/50)}.
\]
By \eqref{smallbranch}, the coefficient of $e$ in parentheses is at least
\[
 \frac{183}{200}-\frac{20}{57}\frac{725029049}{278853120}
                         =\frac{53754367}{19868284800}>0.
\]
Thus every term on the right of \eqref{certificate-start} after $d_0$
is nonnegative. The displayed squares give an explicit certificate for the remaining
quadratic form.

\begin{proposition}\label{conditional-law}
Every isotropic law on $\R^3$ whose radii lie almost surely in a positive
interval of endpoint ratio at most $3$ satisfies $D\ge d_0$.
In the branch $D<1/20$, \eqref{certificate-start}--\eqref{squares}
give the full common-remainder certificate.
\end{proposition}
\begin{proof}
The small branch was proved above. Otherwise $D\ge1/20>d_0$.
For explicit exact comparisons use
$682/305<\sqrt5<161/72$ and $433/250<\sqrt3<97/56$,
each verified by squaring. They yield
\[
 d_0-7/144>\frac{940735364069}{29456638041600000}>0,\qquad
 1/20-d_0>\frac{620701786231}{482895705600000}>0.
\]
No $B=0$ division occurs: it is excluded by $e<1/20$ in the small branch.
\end{proof}


\section{A reserve in the quartic slack}\label{sec-reserve}
% Ported verbatim from paper/top3_joint_moment_bound.tex lines 1086-1103.
% Only cross-reference wording was adapted; no mathematics was altered.
Throughout this section assume $0<r_-\le R\le r_+<\infty$ almost surely,
with $r_+/r_-\le\sqrt8$. Write $g(w)=(1-9w^2)/25$. Then
\[
 \frac{(R-R')^2}{RR'}\le\sqrt8+\frac1{\sqrt8}-2
    =\frac{9\sqrt2}{4}-2\le\frac{32}{27}.
\]
For the last inequality, $\sqrt2<344/243$ follows by squaring, with
$(344/243)^2-2=238/59049>0$. Put
$z_{\rm rad}=3(R-R')^2/(4RR')$, so $0\le z_{\rm rad}\le8/9$.
The actual remainder in \eqref{payment} now has the exact decomposition
\begin{equation}\label{reserve}
 r_{\rm rad}=\frac P9+r_{\rm tr},\qquad
 r_{\rm tr}=\E\!\left[RR'\left\{(8/9-z_{\rm rad})\sigma(w)
                    +z_{\rm rad}(\sigma(w)-g(w))\right\}\right]\ge0.
\end{equation}
This uses both pointwise inequalities of Lemma~\ref{scalar}, including
where $g<0$. It is not a second independent budget or a new copy of $P$.


\subsection{A contact polynomial without a cubic term}
% Ported verbatim from paper/top3_joint_moment_bound.tex lines 1105-1190.
% Only cross-reference wording was adapted; no mathematics was altered.
Using the same interpolation nodes and tensor weights, set
\[
 h_*(t)=(t-1)(t-1/2)(t+3/8),\qquad
 z_*(t)=(t+9/8)h_*(t)
       =t^4-\frac{85}{64}t^2+\frac{15}{128}t+\frac{27}{128}.
\]
The factor $t+9/8$ cancels the cubic term. We need the following elementary
bound, for which we give an explicit proof.
\begin{lemma}\label{contact-square}
For $-1\le t\le1$, $z_*(t)^2\le3\sigma(t)/2$.
\end{lemma}
\begin{proof}
For $0\le t\le1$, cancel common factors away from their zeros to obtain
\[
 \frac{z_*(t)^2}{\sigma(t)}
  =\frac{41503}{4(8352+2560t)}q_*(t),\qquad
 q_*(t)=(1-t)(t+3/8)^2(t+9/8)^2.
\]
The cancelled expressions extend continuously at the zeros. Its derivative is
\[
 q_*'(t)=(t+3/8)(t+9/8)(165/64-t/2-5t^2).
\]
The last factor decreases on $[0,1]$, equals $13/576>0$ at $2/3$,
and equals $-33/256<0$ at $11/16$. Thus the maximum is within
$[2/3,11/16]$. On that interval $(t+3/8)(t+9/8)<2$ and
$q_*'(t)\le13/288$. The interval has length $1/48$, so
\[
 q_*(t)\le\frac{1155625}{995328}+\frac{13}{13824}<\frac65;
\]
the final positive difference is $189163/4976640$.
Consequently $z_*^2/\sigma\le41503/27840<3/2$ on this half interval.
For $t=-x$, $0\le x\le1$, the other factorization gives
\[
 \frac{z_*(-x)^2}{\sigma(-x)}
 =\frac{41503\,p_*(x)^2}{64(928+82x-160x^2)},\qquad
 p_*(x)=(1+x)(x+1/2)(9/8-x)\ge0.
\]
The quadratic denominator factor is at least $850$, by concavity and its
endpoint values. With $y=x-1/2\in[-1/2,1/2]$,
\[
 p_*(x)=15/16+y/16-15y^2/8-y^3
 \le15/16+y/16-11y^2/8
 =\frac{1321}{1408}-\frac{11}{8}(y-1/44)^2<1.
\]
It follows that $z_*^2/\sigma\le41503/54400<3/2$ here.
At every cancelled zero the claimed inequality follows directly or by
continuity. These two intervals cover the domain.
\end{proof}

\begin{remark}\label{rem-theta}
The constant $3/2$ is not sharp: a dense grid scan gives
$\sup_{[0,1]}z_*^2/\sigma=1.200224\ldots$ and
$\sup_{[-1,0]}z_*^2/\sigma=0.615384\ldots$. We keep $3/2$ because it keeps
every subsequent coefficient rational and hand-verifiable. The loss enters
the certificate only through $r_{\rm ct}$ and through the divisor $125$ in
\eqref{Cstar8}; both are quantified in Appendix~\ref{sec-appendix}.
\end{remark}

In particular the following displayed square proves a useful linear payment:
\begin{equation}\label{contact-paid}
 \sigma(t)-\frac{z_*(t)}{125}+\frac3{125000}
   =\left(\sigma(t)-\frac23z_*(t)^2\right)
       +\frac23\left(z_*(t)-\frac3{500}\right)^2\ge0.
\end{equation}
All these functions are bounded on $[-1,1]$. Integrating against the actual
nonnegative weight $RR'$ uses only $\E R<\infty$, not a fourth or eighth
radial moment. The actual contact moment is
\begin{align}\label{Zstar}
 Z_*:=\E[RR'z_*(U\cdot U')]
   &=K^2+\frac7{320}-\frac{15}{128}e
                    -\frac{561}{4480}V+\frac{211}{448}J.
\end{align}
Indeed the $RR'$ moments of $1,w,w^2,w^4$ are respectively
$3-V,1-e,1-J,K^2+6\|Z\|^2/7+(3-V)/5$, all for this same law.
Define
\begin{equation}\label{Cstar8}
 {\cal C}=Z_*/125-3C^2/125000,\qquad
 r_{\rm ct}=P-{\cal C}\ge0.
\end{equation}
The scalar ${\cal C}$ need not itself be nonnegative. Its exact contribution is
\begin{equation}\label{paid-expansion}
 \frac{{\cal C}}{9}=\frac{103}{9000000}+\frac{K^2}{1125}
 -\frac e{9600}+\left(-\frac{187}{1680000}+\frac1{375000}\right)V
                       +\frac{211}{504000}J.
\end{equation}


\subsection{The full joint certificate}
% Ported verbatim from paper/top3_joint_moment_bound.tex lines 1192-1227.
% Only cross-reference wording was adapted; no mathematics was altered.
For $D<1/20$ define the explicitly nonnegative quadratic expression
\begin{align}\label{Q8}
 Q={}&\frac{13}{200}\left(\sqrt V-\frac{937}{273}q\sqrt e\right)^2
       +\frac{19}{50}\left(\sqrt J-\frac{2811}{2128}q\sqrt e\right)^2\notag\\
 &+\left(\frac{183}{200}-\frac{725029049}{278853120}q^2\right)e\notag\\
 &+\frac{148831}{181104000}e
       +\frac{42847433}{5229378000000}V+\frac{65082103}{20917512000}J.
\end{align}
The parenthetical coefficient is positive by \eqref{smallbranch} and the
same exact lower bound following \eqref{squares}. The last coefficients are
obtained from those in \eqref{squares}, respectively by adding
$-1/9600$, $-187/1680000+1/375000$, and $211/504000$.
Substitution of \eqref{reserve}--\eqref{paid-expansion} into \eqref{certificate-start} yields the exact identity
\begin{equation}\label{certificate8}
 D=d_1+Q+{\cal R}_0+\ell r_{\rm reg}
          +\frac{K^2}{1125}+r_{\rm tr}+\frac{r_{\rm ct}}9.
\end{equation}
Every remainder is defined from the same $X$, its independent copy and its
complete Gram matrix.
\begin{proposition}\label{conditional8}
Every isotropic law whose radii lie almost surely in a positive bounded
interval of endpoint ratio at most $\sqrt8$ satisfies $D\ge d_1$.
In the branch $D<1/20$, \eqref{certificate8} retains all displayed remainders.
\end{proposition}
\begin{proof}
In the small branch all terms after $d_1$ are nonnegative. In the other
branch $D\ge1/20>d_1$. To verify the last comparison entirely by rationals,
use $\sqrt5<161/72$ and $\sqrt3>433/250$ in \eqref{constants}; they give
\[
 \frac1{20}-d_1>\frac{1025292188593}{804826176000000}>0.
\]
No division by $B=0$ occurs in the small branch, and the other branch
requires none. The radial hypothesis will be verified for a maximizing law in
Section~\ref{sec-conclusion}, rather than imposed on every isotropic law.
\end{proof}


\section{Transfer to graphs and proof of the main theorem}\label{sec-conclusion}
% Ported verbatim from paper/top3_joint_moment_bound.tex lines 1229-1278.
% Only cross-reference wording was adapted; no mathematics was altered.
We return to a maximizing graphon. Put
\[
 D_*=(3+\sqrt5)/2-2\mathcal M.
\]
By the threshold identity in Lemma~\ref{contact-structure}, its canonical
law $X=F$ satisfies $A+B=2\mathcal M$, so the deficit $D$ of the preceding
sections equals $D_*$.
\begin{lemma}\label{spectral-reserve}
If $W$ is an actual maximizer and $D_*=(3+\sqrt5)/2-2{\cal M}<1/20$,
then $\lambda_3>44/375$, $\lambda_1<15/16$, and the canonical frame
has positive bounded radii of endpoint ratio less than $\sqrt8$.
\end{lemma}
\begin{proof}
The small branch and $\sqrt5>682/305>559/250$ give
\[
 {\cal M}>\frac{3+\sqrt5}{4}-\frac1{40}>\frac{321}{250}.
\]
Lemma~\ref{ranktwo} gives
$\lambda_3={\cal M}-\Sigma_2(W)>321/250-7/6=44/375$.
For $a=\lambda_1$, the Hilbert--Schmidt trace and the constant-function trial give
\[
 \lambda_1^2+\lambda_2^2+\lambda_3^2\le\iint W^2\le\iint W\le\lambda_1,
 \qquad {\cal M}\le a+\sqrt{2a(1-a)},\quad 0<a\le1.
\]
The function on the right is decreasing on $[15/16,1]$; differentiation
on the interior reduces this to $(2a-1)^2>2a(1-a)$ there.
Moreover
\[
 \frac{15}{16}+\sqrt{\frac{15}{128}}<\frac{321}{250},\qquad
 \left(\frac{321}{250}-\frac{15}{16}\right)^2-\frac{15}{128}
       =\frac{11499}{4000000}>0.
\]
Thus $a<15/16$. By the actual ellipsoid identity,
\[
 \sqrt{{\cal M}/\lambda_1}\le R\le\sqrt{{\cal M}/\lambda_3},\qquad
 \frac{\esssup R}{\essinf R}
 \le\sqrt{\lambda_1/\lambda_3}<\sqrt{5625/704}<\sqrt8,
\]
where $8-5625/704=7/704$. No simplicity assumption was made on
the positive eigenvalues; the argument uses the whole orthonormal frame.
\end{proof}

\begin{remark}\label{rem-radii}
Lemma~\ref{radialpayment} is stated with endpoint ratio $3$, whereas the
present lemma supplies $\sqrt8<3$ for an actual maximizer; the two
hypotheses are therefore compatible, with slack. Using $\sqrt8$ in
Lemma~\ref{radialpayment} would replace the factor $4/3$ in \eqref{payment}
by $32/27$, a $12.5\%$ gain on $P$, which is not needed here.
\end{remark}


\subsection{Positive-mass self-pairs}
% Ported verbatim from paper/top3_joint_moment_bound.tex lines 1280-1346.
% Only cross-reference wording was adapted; no mathematics was altered.
By Theorem~\ref{structure-main}, write the canonical law as
$\sum_{i=1}^m p_i\delta_{v_i}$ with $m\le24$ and all $p_i>0$.
Set $R_i=|v_i|$, $u_i=v_i/R_i$, $w_{ij}=u_i\cdot u_j$, and
\[
 \kappa=\frac3{125000},\qquad
 h(t)=\sigma(t)-\frac{z_*(t)}{125}+\kappa,
 \qquad \chi=\sum_i p_i^2R_i^2.
\]
Equation~\eqref{contact-paid} gives $h\ge0$ and $h(1)=\kappa$.
The same contact remainder in~\eqref{Cstar8} therefore has the exact split
\begin{equation}\label{atom-split}
 r_{\rm ct}=\kappa\chi+r_{\rm off},\qquad
 r_{\rm off}=\sum_{i\ne j}p_ip_jR_iR_jh(w_{ij})\ge0.
\end{equation}
These self-pairs have positive product mass. They are not values assigned
to the null diagonal $x=y$ of the underlying nonatomic space.
The estimate $\chi\ge C^2/m$ is the usual Cauchy--Schwarz
cardinality bound. For related weighted-energy arguments,
see~\cite[Theorem 2.1 and the subsequent discussion]{Weighted}.

In the branch $D_*<1/20$, let
\[
 \widehat Q=Q-\frac{V}{9000000}.
\]
It is nonnegative: in~\eqref{Q8} only the last separate $V$ coefficient
changes, to
\[
 \frac{42847433}{5229378000000}-\frac1{9000000}
       =\frac{14088797}{1743126000000}>0.
\]
\begin{proposition}\label{final-certificate}
For every maximizing graphon with $D_*<1/20$,
\begin{align}\label{certificate-final}
 D_*={}&d+\widehat Q+{\cal R}_0+\ell r_{\rm reg}
       +\frac{K^2}{1125}+r_{\rm tr}\notag\\
 &+\frac{\kappa}{9}\left(\chi-\frac{C^2}{24}\right)
       +\frac{r_{\rm off}}9.
\end{align}
Every term after $d$ is nonnegative. Consequently $\mathcal M\le c$.
\end{proposition}
\begin{proof}
Lemma~\ref{spectral-reserve} verifies the radial hypothesis of
Proposition~\ref{conditional8}. Substitute~\eqref{atom-split} once into
\eqref{certificate8}. Since $C^2=3-V$,
\[
 \frac{\kappa\chi}{9}
 =\frac1{3000000}-\frac{V}{9000000}
       +\frac{\kappa}{9}\left(\chi-\frac{C^2}{24}\right).
\]
This gives~\eqref{certificate-final}. Moreover
\[
 \chi-\frac{C^2}{24}
 =\sum_i\left(p_iR_i-\frac Cm\right)^2
       +C^2\left(\frac1m-\frac1{24}\right)\ge0.
\]
No lower bound on an individual mass is needed. The complementary branch
also gives $D_*\ge d$, since the rational estimates used above imply
\[
 \frac1{20}-d>
       \frac{1025023913201}{804826176000000}>0.
\]
Thus $D_*\ge d$ in both branches, or equivalently $\mathcal M\le c$.
\end{proof}
The atomic term changes the coefficient by $1/6000000$ relative to
$(3+\sqrt5)/4-d_1/2$. Most of the improvement in Theorem~\ref{main}
comes from the preceding joint moment estimate, not from this last term.


\subsection{All isotropic laws and the exact finite shift}
% Ported verbatim from paper/top3_joint_moment_bound.tex lines 1348-1384.
% Only cross-reference wording was adapted; no mathematics was altered.
\begin{proof}[Proof of Theorem~\ref{main}]
Any Borel law on $\R^3$ can be realized on a standard nonatomic probability
space, by adjoining an independent uniform variable if necessary.
For its isotropic $X$, the coordinates are an orthonormal three-frame and
$W_X=\mathbf1_{\{X\cdot X'>0\}}$ is a measurable symmetric binary kernel.
Therefore
\[
 \frac{\E|X\cdot X'|+|\E X|^2}{2}
   =\iint W_X X\cdot X'
   \le\Sigma_3(W_X)\le\mathcal M\le c.
\]
This proves~\eqref{law-bound}, including zero vectors and unbounded radii.
The radial restriction was used only for a maximizing canonical law.

For a finite graph, the exact-rank Ky Fan formula reads
\[
 \Lambda_3(G)=\max_{Q=Q^T=Q^2,\ \tr Q=3}\tr(A_GQ).
\]
Factor $Q=YY^T$ with orthonormal columns, let $y_i$ be its rows, and let
$X$ be uniform on $\sqrt n\,y_1,\ldots,\sqrt n\,y_n$. Then
\[
 \E XX^T=I_3,\qquad
 \E|X\cdot X'|=\|Q\|_1/n,\qquad
 |\E X|^2=\boldsymbol1^TQ\boldsymbol1/n,
\]
where $\|Q\|_1$ denotes the entrywise absolute sum.
Since $Q_{ii}\ge0$ and $\sum_iQ_{ii}=3$, the zero adjacency diagonal gives
\[
 \tr(A_GQ)\le\sum_{i\ne j}(Q_{ij})_+
 =\frac{\|Q\|_1+\boldsymbol1^TQ\boldsymbol1}{2}-3
 \le cn-3.
\]
Taking the maximum proves~\eqref{main-bound}. No positivity assumption on
$\lambda_3(G)$ is used; repeated eigenvalues, zero rows and $n=3$ cause
no exception.
\end{proof}


\section{Exact scope of the bound}\label{sec-weighted}

The transfer argument above uses two properties of $A_G$ only: its diagonal is
zero, and its off-diagonal entries lie in $[0,1]$. Both are shared by a larger
class, and the bound therefore holds there without further work. The scope of such a statement is easy to mis-state, so we record it
precisely, and we record at the same time that no novelty is claimed for it:
the same observation underlies the proof of~\cite[Proposition 6.1]{STW}, and
the matrix setting is explicit in the title and content of~\cite{Mohar}.

\begin{proposition}[Weighted form]\label{weighted-form}
Let $n\ge3$ and let $A\in\R^{n\times n}$ be symmetric with $A_{ii}=0$ and
$0\le A_{ij}\le1$ for $i\ne j$. Then
\[
 \lambda_1(A)+\lambda_2(A)+\lambda_3(A)\le cn-3,
\]
with the constant $c$ of~\eqref{main-constant}.
\end{proposition}
\begin{proof}
Apply Lemma~\ref{bridge} with $r=3$. The isotropic law produced there is
built from the rows of a rank-three projection and satisfies no radial
hypothesis, so Theorem~\ref{main} applies to it and gives $\Gamma_3\le c$.
The constant term is $-r=-3$ and is unchanged.
\end{proof}
\begin{remark}\label{weighted-scope}
Proposition~\ref{weighted-form} covers $[0,1]$-weighted loopless graphs, hence
also every graph with edge weights in $[0,1]$ arising from a discretised
graphon. It fails for weights above $1$ and for multiple edges: if $A=2A_G$
then $\Lambda_3(A)=2\Lambda_3(G)$, which exceeds $cn-3$ for large $n$
whenever $\Lambda_3(G)/n> c/2$. The numerical tests of
Section~\ref{sec-num-weighted} are consistent with this.
\end{remark}

Weights cannot improve the asymptotic constant either, so the extremal
problem does not become easier by relaxation.
\begin{proposition}[Weights give no asymptotic gain]\label{no-gain}
Let $\mathcal M_{\rm w}(n)$ be the supremum of $\Lambda_3(A)$ over symmetric
$A\in\R^{n\times n}$ with $A_{ii}=0$ and $0\le A_{ij}\le1$. Then
\[
 \frac{\mathcal M_{\rm w}(n)}n\ \longrightarrow\ \mathcal M
 \qquad(n\to\infty),
\]
the same limit as for unweighted graphs.
\end{proposition}
\begin{proof}
Every such $A$ defines a step graphon $W_A$ with $T_{W_A}=A/n$, so by
\eqref{wg-spectrum} $\Lambda_3(A)\le n\Sigma_3(W_A)\le n\mathcal M$; hence
$\limsup_n\mathcal M_{\rm w}(n)/n\le\mathcal M$. For the reverse inequality,
$\mathcal M$ is attained by some graphon $W$ by Lemma~\ref{attainment}. By the
compactness and density statements of Section~\ref{sec-graphons} there are
graphs $G_m$ with $\delta_\square(W_{G_m},W)\to0$, and $\Sigma_3$ is
continuous in $\delta_\square$ because
$|\Sigma_3(W_1)-\Sigma_3(W_2)|\le3\|T_D\|_{2\to2}\le3(4\|D\|_\square)^{1/3}$
by~\eqref{cut-eigen}. Writing $n_m=|V(G_m)|$, the identity
$\Lambda_3(G_m)=n_m\Sigma_3(W_{G_m})+\sum_{i\le3}\min(\lambda_i(G_m),0)$
and the bound $\lambda_i(G_m)\ge-(n_m-1)$ give
\[
 \frac{\mathcal M_{\rm w}(n_m)}{n_m}
   \ge\frac{\Lambda_3(G_m)}{n_m}
   \ge\Sigma_3(W_{G_m})-3\,\frac{n_m-1}{n_m}
   \longrightarrow\mathcal M .
\]
Replacing $G_m$ by a blow-up if necessary, which does not change its graphon,
we may assume $n_m\to\infty$. Since $\mathcal M_{\rm w}(n)/n\le\mathcal M$
for every $n$, the limit exists and equals $\mathcal M$.
\end{proof}
\begin{remark}\label{padding-note}
The padding term is nonnegative: by~\eqref{wg-spectrum},
\[
 n\Sigma_3(W_G)-\Lambda_3(G)=-\sum_{i\le3}\min(\lambda_i(G),0)\ge0 .
\]
It is therefore harmless in the direction used above, where only a lower bound
for $\mathcal M_{\rm w}(n)$ is needed.
\end{remark}
